% Part: computability
% Chapter: computability-theory
% Section: def-functions-self-reference

\documentclass[../../../include/open-logic-section]{subfiles}

\begin{document}

\olfileid{cmp}{thy}{slf}
\olsection{સ્વ-સંદર્ભ વડે વિધેયોની વ્યાખ્યા}

વિધેયોને તેમના પોતાના જ આધારે વ્યાખ્યાયિત કરી શકવાની
ક્ષમતા સામાન્ય રીતે ઉપયોગી છે. દાખલા તરીકે, સંગણનીય
વિધેયો $k$, $l$ અને~$m$ આપેલાં હોય, તો નિશ્ચિત-બિંદુનું
ઉપપ્રમેય કહે છે કે દરેક~$y$ માટે નીચેનું સમીકરણ સંતોષતું
આંશિક સંગણનીય વિધેય~$f$ છે:
\[
f(y) \simeq
\begin{cases}
k(y) & \text{જો $l(y) = 0$} \\
f(m(y)) & \text{અન્યથા.}
\end{cases}
\]
વધુ ચોક્કસ રીતે, $f$ મેળવવા પહેલાં
\[
g(x,y) \simeq
\begin{cases}
k(y) & \text{જો $l(y) = 0$} \\
\cfind{x}(m(y)) & \text{અન્યથા}
\end{cases}
\]
રાખીએ અને પછી નિશ્ચિત-બિંદુના ઉપપ્રમેયથી એવો સૂચક~$e$
મેળવીએ કે $\cfind{e}(y) = g(e,y)$.

ચોક્કસ ઉદાહરણ તરીકે, ``મહત્તમ સમાપવર્તક'' વિધેય
$\fn{gcd}(u,v)$ આ રીતે વ્યાખ્યાયિત કરી શકાય:
\[
\fn{gcd}(u,v) \simeq
\begin{cases}
v & \text{જો $u = 0$} \\
\fn{gcd}(\fn{mod}(v, u), u) & \text{અન્યથા}
\end{cases}
\]
અહીં $\fn{mod}(v, u)$, $v$ને~$u$ વડે ભાગતાં મળતી
શેષ સંખ્યા દર્શાવે છે. નિશ્ચિત-બિંદુના ઉપપ્રમેયથી
$\fn{gcd}$ આંશિક સંગણનીય છે. (હકીકતમાં $y$ જોડી
$\tuple{u, v}$નો કોડ હોય એમ લઈને તેને ઉપરના સ્વરૂપમાં
લાવી શકાય.) ત્યાર પછી~$u$ પર અનુમાનથી બતાવી શકાય કે
$\fn{gcd}$ ખરેખર સંપૂર્ણ છે.

અલબત્ત, હમણાં ચર્ચેલાં ઉદાહરણો કરતાં ઘણાં વધુ જટિલ
સ્વ-સંદર્ભી વ્યાખ્યાઓ આપી શકાય છે. મોટાભાગની પ્રોગ્રામિંગ
ભાષાઓ કોઈ ને કોઈ રીતે વિધેયને તેના પોતાના આધારે વ્યાખ્યાયિત
કરવાની છૂટ આપે છે. આ, વિધેયની ગણતરી કરતી કલનવિધિના
\emph{સૂચક}ના આધારે વિધેય વ્યાખ્યાયિત કરવાની ક્ષમતા કરતાં
થોડું ઓછું અદ્ભુત છે; નિશ્ચિત-બિંદુનું પ્રમેય તેની પૂર્ણ
સામાન્યતામાં પછીની ક્ષમતા આપે છે.

\end{document}
